\documentclass[11pt,a4paper]{article}
\usepackage[T1]{fontenc}
\usepackage[utf8]{inputenc}
\usepackage{lmodern}
\usepackage[margin=2.3cm]{geometry}
\usepackage{amsmath,amssymb}
\usepackage{booktabs}
\usepackage{graphicx}
\usepackage{xcolor}
\usepackage[colorlinks=true,linkcolor=blue!50!black,urlcolor=blue!50!black]{hyperref}
\usepackage{caption}
\usepackage{float}
\captionsetup{font=small,labelfont=bf}

\newcommand{\pass}{\textcolor{green!45!black}{\textbf{PASS}}}
\newcommand{\fail}{\textcolor{red!70!black}{\textbf{FAIL}}}
\newcommand{\code}[1]{\texttt{\hyphenchar\font=`\-#1}}
\sloppy
\emergencystretch=2em

\title{Pair functions in a hadron resonance gas at rest:\\ results of milestone M6
(single-particle currents)}
\author{THERMINATOR 2 (hboltz fork), \code{src/work/hrg\_pair\_density\_spec.md}, Sec.~9}
\date{26 October 2026}

\begin{document}
\maketitle

\begin{abstract}
\noindent
This note reports milestone M6 of Part~II of the spec: the single-particle currents
$S_c^\mu=\mathcal A_c\,p^\mu+\mathcal B_c\,\epsilon u^\mu$ of all 381 species of the THERMINATOR
\code{ParticleDB}. They are built by a FastReso-type recursion over the parsed decay channels, in
decreasing mass. All tests pass for both weak-decay settings: T6.1 (one step from thermal parents
equals the closed-form two-body $F_1$, $F_1+F_u$ to $1.5\times10^{-11}$), T6.2 (densities equal the
counting to $1.5\times10^{-8}$, cut-independent to $3\times10^{-10}$), T6.3 (positivity, timelike,
finite at $k\to0$) and T6.4 (event MC, $2\times10^4$ events: species densities and momentum spectra).
M7 has not been started.
\end{abstract}

\section{Implementation}
Code: new section ``Single-particle currents'' in \code{src/work/PairFunctions.h}, new program
\code{src/work/ex\_pair\_functions\_currents.exe}, and a \code{--spectra} flag in
\code{ex\_pair\_functions\_mc.exe}.

\paragraph{Representation.} For every species, the thermal part
$\mathcal A^{\rm th}_c=s_c\,g_cf_c(\epsilon)/(2\pi)^3$, $\mathcal B^{\rm th}_c=0$, is kept analytic.
Here $s_c=1$ (\code{--norm quad}) or $s_c=$ \code{GetMultiplicity()}$/(Vn_c^{\rm quad})$
(\code{--norm mc}, the multiplicities the event MC samples), and $s_c=0$ for photons. The feed-down
part ($\mathcal A_c$ and $\mathcal A_c+\mathcal B_c$) is stored on a momentum grid. The grid has the
panels of the M1 table (\code{ThermalPanels}: $0.5T$ up to $10T$, $T$ up to $40T$, $5T$ up to the
cut, default $100T$), with 16 Gauss--Legendre nodes per panel (992 nodes). Between nodes the feed-down
part is evaluated by Lagrange interpolation on the panel's nodes. It is zero beyond the cut.

\paragraph{Recursion.} There are 3597 feed slots with weak decays on (3282 two-body, 315
three-body), 3549 with weak decays off. They give every daughter slot of every channel of every
decaying parent, with the BR as stored after parsing. Species are processed in decreasing mass, so
the parents of a species are complete when it is reached. Weak-stable species (\code{--weak off}) and
species with BR sum below 1 are handled as in the table. A three-body slot is the two-body rule
averaged over the effective partner mass $\tilde m\in[m_c+m_d,m_a-m_b]$, with weight
$p^*(m_a;m_b,\tilde m)\,p^*(\tilde m;m_c,m_d)$ (32 nodes, normalized).

\paragraph{Deviation from the spec: integration variable.} The two-body rule is integrated in the
parent energy $E$ instead of $w$. This is the same integral, since $E(w)$ is linear in $w$:
\begin{align*}
\mathcal A_b&\mathrel{+}=\frac{Bm_a}{2p^*k^3}\int_{E^-}^{E^+}\!dE\,(\epsilon E-m_aE_b^*)\,\mathcal A_a(E),&
(\mathcal A_b+\mathcal B_b)&\mathrel{+}=\frac{Bm_a}{2p^*k\epsilon}\int_{E^-}^{E^+}\!dE\,E\,(\mathcal A_a+\mathcal B_a)(E),
\end{align*}
with $E^-=m_a(E^{*2}+k^2)/(\epsilon E^*+kp^*)$ (no $1/m_b^2$) and the upper limit clipped at the
parent's grid cut. The panels start at $E^-$ and follow the thermal scale
($E^-+T\{0,\tfrac12,1,2,4,\dots,64,100\}$). A fixed rule in $w$ fails at large $k$: for pions from
heavy parents the window $E^+-E^-$ is $\gg T$. For narrow windows ($\epsilon h<m_aE^*$, with
$h=m_akp^*/m_b^2$ the half-width) the factor is written as $\epsilon y+m_aE^*k^2/m_b^2$, with
$y=E-E_c$ on the exactly symmetric window $[E^-,E^-+2h]$. Without this, the $1/k^3$ cancellation
cost three digits at the smallest node: the first run gave $2.7\times10^{-9}$ in T6.1, and the
feed part changed by a factor 2 between $k=10^{-6}$ and $10^{-8}$.

\paragraph{Photons} (\code{gam000zer}, $m=0.001$~GeV in \code{particles.data}) are not primordial
but are produced (e.g.\ $\Sigma^0\to\Lambda\gamma$). Their current is built like any other and is
included in T6.2 and T6.4. It is reported separately in T6.3.

\paragraph{Run time.} 40~s per build of all currents on 8 threads. A full run with the cut check
takes about 75~s.

\section{Results}
Commands (in \code{build/t2/}, which holds the \code{fmultiplicity\_E33B55BA.txt} cache;
\code{THERMINATOR\_SHARE} set):
\begin{verbatim}
ex_pair_functions_currents.exe --weak on|off \
    --output m6/currents_weak{on,off}_quad.root
ex_pair_functions_currents.exe --weak on --stat boltzmann --cutcheck 0 ...
ex_pair_functions_mc.exe --spectra --weak on|off --events 5000 --seed 601..604   (hadd)
ex_pair_functions_currents.exe --weak on|off --norm mc --cutcheck 0 \
    --mc m6/spectra_weak{on,off}.root --output m6/currents_weak{on,off}_mc.root
\end{verbatim}

\begin{table}[H]
\centering\footnotesize
\begin{tabular}{llll}
\toprule
test & quantity & worst deviation (weak on / off) & \\
\midrule
T6.1 & $\mathcal A_c$ vs closed-form $F_1$, $2.3\times10^6$ (slot, node) pairs & $6.2\times10^{-13}$ / $6.2\times10^{-13}$ & \pass\\
T6.1 & $\mathcal A_c+\mathcal B_c$ vs $F_1+F_u$ & $1.5\times10^{-11}$ / $1.5\times10^{-11}$ & \pass\\
T6.2 & $n_c^{\rm grid}$ vs counting, 381 species & $1.5\times10^{-8}$ ($\gamma$) / same & \pass\\
T6.2 & cut $40T$ vs $100T$, $\max|\Delta n/n|$ & $2.8\times10^{-10}$ ($J/\psi$) / same & \pass\\
T6.3 & 376\,960 hadron nodes: non-finite, $\mathcal A+\mathcal B\le0$, spacelike & 0, 0, 0 / 0, 0, 0 & \pass\\
T6.3 & $k=10^{-6},10^{-8}$: finite, $\mathcal A+\mathcal B>0$; rel.\ change & 67/67; $1.3\times10^{-8}$ & \pass\\
T6.4 & species densities, pulls (381 species) & rms 1.01, max 2.8 / rms 0.97, max 3.7 & \pass\\
T6.4 & $dN/dp$, 8 species, $\chi^2$ & min $p=0.035$ ($\Lambda$) / $0.045$ ($\rho^0$) & \pass\\
\bottomrule
\end{tabular}
\caption{M6 tests. The \code{--stat boltzmann} run gives the same verdicts and deviations of the
same size.}
\end{table}

\paragraph{T6.1.} The closed form is LCS Eq.~(twobodyF), with $f_a$ expanded as
$\sum_nc_ne^{-nE/T}$ (BE: $c_n=1$, FD: $c_n=(-1)^{n-1}$) and each term integrated analytically
(with a series for $\sinh x-x\cosh x$ when $x=nh/T<0.5$). Nodes where the parent grid cut enters
($E^->E_{\max}-50T$, i.e.\ a truncation above $e^{-50}$; 28\% of the nodes, all in the far tail)
are excluded from the comparison. The worst cases are at the smallest node ($k=4.1\times10^{-4}$
GeV).

\paragraph{T6.2.} Spec Sec.~9.5 counting in the same order:
$n_b=n_b^{\rm th}+\sum_an_a\sum_{\rm ch}B\,n_b^{\rm ch}$. Selected densities (weak on, $n_c/n_c^{\rm th}$):
$\pi^+$ 3.11, $K^+$ 1.83, $p$ 4.17, $\Lambda$ 5.29, $\rho^0$ 1.38, $\omega$ 1.36, $K^{*0}$ 1.22,
$\Delta^{++}$ 1.29.

\paragraph{T6.3.} $\mathcal A_c$ turns out to be positive on every node, although the spec does not
require it. The smallest timelike margin $(\epsilon(\mathcal A+\mathcal B)-k|\mathcal A|)/
(\epsilon(\mathcal A+\mathcal B))$ is $3.8\times10^{-5}$ ($\pi^0$ at large $k$, where
$\epsilon\to k$). Only 67 species receive feed-down, and only these enter the $k\to0$ check.
Fig.~\ref{fig:cur} shows that $\mathcal B_c/(\mathcal A_c+\mathcal B_c)$ reaches 0.77 for slow
$\pi^+$: the $u^\mu$ part of the current matters for M7.

\paragraph{T6.4.} The event MC fills, for all particles in the list (decayed or not), the per-event
count and count$^2$ of each species, and $dN/dp$ (10~MeV bins, 0--3~GeV) for $\pi^+$, $K^+$, $p$,
$\rho^0$, $K^{*0}$, $\omega$, $\Lambda$, $\Delta^{++}$. Weak on: per event, $\pi^+$ $3834.33\pm0.46$
(MC) vs 3834.08 (currents), $\Lambda$ $161.65\pm0.09$ vs 161.70, $p$ $322.38\pm0.13$ vs 322.41.
Weak off: $\pi^+$ $3262.82\pm0.41$ vs 3262.97, $\Lambda$ $115.68\pm0.08$ vs 115.64. Spectra:
Fig.~\ref{fig:spec}; the $\chi^2$ test uses bins with at least 5 expected counts. The weak-on and
weak-off samples use the same seeds (601--604), so they are not statistically independent.

\begin{figure}[H]
\centering
\includegraphics[width=\textwidth]{fig/spectra_weakon.png}
\caption{T6.4, weak on: $dN/dp$ of all particles of each species (decayed or not), $2\times10^4$
events, MC (points) vs $N_{\rm ev}V\int 4\pi p^2(\mathcal A_c+\mathcal B_c)\,dp$ (red).
Weak off: \code{fig/spectra\_weakoff.png}.}
\label{fig:spec}
\end{figure}

\begin{figure}[H]
\centering
\includegraphics[width=0.6\textwidth]{fig/pulls.png}
\caption{T6.4: pulls of the species densities, 381 species, both weak settings.}
\end{figure}

\begin{figure}[H]
\centering
\includegraphics[width=\textwidth]{fig/currents.png}
\caption{Currents (weak on, \code{--norm quad}): $\mathcal A_c+\mathcal B_c$ and
$\mathcal B_c/(\mathcal A_c+\mathcal B_c)$, which is zero for a purely thermal species.}
\label{fig:cur}
\end{figure}

\section{Outputs and status}
\code{build/t2/m6/}: \code{currents\_weak\{on,off\}\_\{quad,mc\}.root},
\code{currents\_weakon\_boltz.root} (TTree \code{currents}: name, pdg, mass, $n^{\rm th}$,
$n^{\rm count}$, $n^{\rm grid}$, $k$, $w$, $\mathcal A$, $\mathcal A+\mathcal B$ (total); with
\code{--mc}, also \code{hSpecMC\_<pdg>}, \code{hSpecPred\_<pdg>}, \code{hPullSpecies}), the merged MC
\code{spectra\_weak\{on,off\}.root} and the logs (copied here). M6: all tests \pass. Next: M7
(L1 table) uses \code{BuildCurrents} in place of the thermal $f_a$.

\end{document}
